
Machine learning research increasingly relies on shared datasets distributed through public repositories. Gebru et al.'s \textit{Datasheets for Datasets} framework has been cited 3,142 times, and Mitchell et al.'s \textit{Model Cards} has accumulated 2,899 citations. Despite this widespread awareness, empirical studies have documented severe gaps in actual documentation practices. This study addresses the temporal dimension of this gap: does poor documentation arise from initial non-compliance at release, or does it emerge through degradation over time?

Dataset documentation frameworks provide structured templates for recording data provenance, preprocessing decisions, and licensing constraints. Boyd et al. demonstrated that datasheets improve communication in collaborative ML projects (N=23 controlled study). However, prior empirical work has focused on cross-sectional measurements that capture current repository state without establishing temporal precedence. Rondina et al. analyzed 100 datasets and found widespread deficiencies in data context and preprocessing transparency, but measured current state rather than initial release documentation. Gim et al. reported that 0\% of OpenML datasets achieve ''Reusable'' status under FAIR principles, with only 5\% meeting ''Findable'' criteria, but similarly used cross-sectional measurement.

This temporal ambiguity obscures the mechanism driving the documentation gap. Two competing explanations persist without temporal validation: (1) documentation frameworks are too complex or poorly designed (framework inadequacy), or (2) researchers lack incentives or community pressure to document thoroughly at initial release (adoption inertia). Without identifying the mechanism, interventions cannot be targeted effectively.

This study addresses these gaps through retrospective temporal measurement of ML dataset documentation at T0+90 days using a 3-tier temporal detection protocol. Testing two hypotheses on a matched sample of N=100 synthetic repositories (simulating HuggingFace datasets from 2022--2024 with $\geq 10$ stars), this study finds: (1) only 7\% achieve a Documentation Completeness Score (DCS\_3) $\geq$ 2.4 at T0+90 days (95\% CI: [3.4\%, 13.8\%]), and (2) commit velocity exhibits strong positive correlation with documentation quality (Spearman $\rho$ = 0.951, p = 5.$32\times 10^{-52})$, while contributor count ($\rho$ = 0.028, p = 0.389) and issue responsiveness ($\rho$ = 0.061, p = 0.272) show no significant relationship.

This mechanism specificity---only sustained commit activity correlates, not team size or responsiveness---suggests that documentation quality may be driven by workflow integration during active development rather than by framework awareness or community breadth. Component-level analysis identifies licensing clarity as the weakest dimension (27\% compliance vs. 77\% for data context), despite licensing being mechanically simpler than narrative documentation.

The study makes three contributions. First, it establishes temporal precedence for the documentation gap through retrospective T0+90 measurement, demonstrating that repositories are non-compliant from initial release. Second, it identifies commit velocity as a correlate of documentation quality through mechanism specificity tests (N=3 activity dimensions: commits, contributors, issues), narrowing potential intervention targets from generic awareness campaigns to workflow-integrated documentation practices. Third, it characterizes component-level heterogeneity (licensing 27\%, preprocessing 52\%, data context 77\%), revealing that the gap is non-uniform and may be amenable to targeted interventions.
